\subsection{Finite-cover coordinates for the cubic restriction}

In this subsection \(B_g\) is the full adelic normalized Alexander
unit, evaluated through a finite character at each coefficient prime.
In particular its evaluation at a cube root in characteristic
\(\ell\ne3\) is not obtained by substituting into a dyadic formal
power series. Put \(G=\widehat{\mathrm{GT}}\),
\(H=G_{\mathbf Q}\), \(N=\ker(\lambda,\rho_2)\), and
\(D=N_{\rm rad}\).
For every prime \(\ell\ge5\), choose a primitive cube root
\(\omega\in\overline{\mathbf F}_\ell\), and write
\[
 n_\ell=\frac{\ell-\chi_{-3}(\ell)}3,\qquad
 \omega^{c_{\ell,3}(g)}
       =B_{g,\ell}(\omega,\omega)^{n_\ell}\quad(g\in N).
\]
The exponent \(c_{\ell,3}(g)\) belongs to \(\mathbf F_3\).
The notation \(\mathbf F_3(1)\) below means the additive group
\(\mathbf F_3\) with \(G\)-action through \(\lambda\bmod3\).

\begin{proposition}\label{prop:cubic-finite-relative-cocycle}
The map \(c_{\ell,3}:N\to\mathbf F_3(1)\) is a continuous
homomorphism, independent of the choice of \(\omega\), and satisfies
\[
 c_{\ell,3}(aga^{-1})=\lambda_a c_{\ell,3}(g).
\]
The homomorphism
\[
 \Xi_{\ell,3}(g)=c_{\ell,3}(g)+
                  \ell\,\rho_\ell^I(g)\bmod3
                       \qquad(g\in N)
\]
vanishes on \(N\cap H\). It extends uniquely to a continuous
crossed homomorphism \(G\to\mathbf F_3(1)\) that is zero on \(H\).
On \(D\) its value is \(c_{\ell,3}\).
\end{proposition}
\begin{proof}
The multiplier-one beta law is multiplicative. Its unitary identity
gives
\[
 B_{g,\ell}(\omega,\omega)
 B_{g,\ell}(\omega^{-1},\omega^{-1})=1.
\]
The first value lies in the residue norm-one torus, of order
\(\ell-\chi_{-3}(\ell)=3n_\ell\). Its \(n_\ell\)-th power
therefore belongs to \(\mu_3\). Replacing \(\omega\) by
\(\omega^{-1}\) inverts both the left side of the defining equation
and its chosen generator, leaving the exponent \(c_{\ell,3}\)
unchanged. This proves well-definedness, continuity and additivity.
The semilinear beta law on \(N\) gives
\[
 B_{aga^{-1}}(X,Y)
       =B_g(X^{\lambda_a},Y^{\lambda_a}).
\]
At a cube root this substitutes either \(\omega\) or
\(\omega^{-1}\), and proves the displayed equivariance.
The Kummer coordinate has the same equivariance.

For \(\sigma\in N\cap H\), use the full arithmetic Anderson
formula in the beta form of
\cite[Remark~2.4, p.~101]{Furusho2003}, with \(s=1/3\) and
the multiplication index equal to \(\ell\). Reduce at the
coefficient prime \(\ell\). The denominator is one, since every
\(\ell\)-th root reduces to one and \(B_\sigma(1,\omega)=1\).
In the numerator the factors with \(k=0,2\bmod3\) are one:
these are the boundary values
\(B_\sigma(\omega,1)\) and \(B_\sigma(\omega,\omega^{-1})\).
The factors with \(k=1\bmod3\) equal
\(B_\sigma(\omega,\omega)\). There are \(n_\ell\) such factors
among \(0\le k<\ell\): the count is \(r\) when
\(\ell=3r+1\), and \(r+1\) when \(\ell=3r+2\).
The right side is \(\omega^{-\ell\rho_\ell^I(\sigma)}\).
Consequently
\[
 c_{\ell,3}(\sigma)=-\ell\rho_\ell^I(\sigma)\pmod3,
\]
as required. This step uses the full formula, including its
constant root of unity; its logarithmic derivative is insufficient.

Use \(G=NH\) and define
\(\Xi_{\ell,3}(nh)=\Xi_{\ell,3}(n)\).
Vanishing on \(N\cap H\) makes this well-defined.
For \(n_i\in N,h_i\in H\), equivariance gives
\[
 \Xi_{\ell,3}(n_1h_1n_2h_2)
   =\Xi_{\ell,3}(n_1)+\lambda_{h_1}\Xi_{\ell,3}(n_2),
\]
which is exactly the crossed law, since \(\lambda_{n_1}=1\).
The continuous surjection \(N\times H\to G\) from a compact
space proves continuity of the descended map. The factorization
\(G=NH\) also proves uniqueness.
\end{proof}

The beta scalar here is detected by a finite cover.
The degree-three cover
\[
                 v^3=t(t-1)
\]
has monodromy \(x\mapsto1,\ y\mapsto1\) in
\(\mathbf Z/3\). Its compactification has equation
\[
                Y^2=4v^3+1,\qquad Y=2t-1,
\]
and genus one: it is fully ramified at the three deleted points,
so Riemann--Hurwitz gives \(2g-2=-6+3(3-1)=0\).
The two nontrivial deck-character eigenspaces in its homology
modulo \(\ell\) have dimension one. The punctured cover has
rank-four first homology; its two-dimensional boundary subspace
is deck-trivial because each of its three punctures is fixed.
Thus the two primitive eigenspaces survive compactification.
For \(g\in N\), the standard automorphism of the punctured line
preserves this cyclic subgroup and the boundary classes. On the
Alexander module its commutator generator is multiplied by
\(B_g(X,Y)\). Specializing \(X=Y=\omega\) therefore gives precisely
the scalar \(B_{g,\ell}(\omega,\omega)\) on the corresponding
elliptic eigenspace. This argument only restricts an existing
GT automorphism. It does not assert that every norm-one scalar
lifts to a GT automorphism.

Finally apply the established integral cubic-frame comparison
and the already eliminated local conditions at two and three.
On \(D\), the kernel of the cubic radical--CM defect has the
explicit description
\[
 \ker(d_{(3)}|_D)
       =\bigcap_{\ell\ge5}\ker(c_{\ell,3}|_D).
\]
Indeed the kernel of the \(n_\ell\)-th power on the cyclic residue
torus is its cube subgroup, and cubing is an automorphism on the
pro-\(\ell\) reduction kernel for \(\ell\ne3\).
Consequently, for the established normalization \(G=R_DH\),
\[
 R_D\cap\ker d_{(3)}
       =R_D\cap\bigcap_{\ell\ge5}\ker\Xi_{\ell,3}.
\]
This makes the permitted cubic restriction in the maximal
pro-two relative quotient a concrete intersection of finite-cover
kernels. It gives no nonzero value of a relative cocycle and does
not assert that the intersection is a proper subgroup of \(R_D\).
